% answers_short.tex —— 2006 年·四川卷·理科 · 简短答案速查
% 由 tools/make_short_answers.py 自 answers.tex 的 \gksolhead/\gkanswer 提取，
% 并与 2006/2006_数学真题_四川理_答案要点.md 逐题比对后生成，勿手改（改源 MD 后重跑）。

\begin{gknotice}
\textbf{使用说明}：本节只给最终答案便于快速核对，完整推导见解析版。答案为编写者独立推导，\textbf{非官方标准答案}。
\end{gknotice}

\gkpart{选择题}
\begin{center}\begin{tabular}{ccccccccccccc}
\toprule
题号 & 1 & 2 & 3 & 4 & 5 & 6 & 7 & 8 & 9 & 10 & 11 & 12 \\
\midrule
答案 & C & D & D & B & D & B & A & C & A & C & A & B \\
\bottomrule
\end{tabular}\end{center}

\gkpart{填空题}
\begin{itemize}[leftmargin=2.6em]
  \item[\textbf{13.}] $\arctan\sqrt2$
  \item[\textbf{14.}] $\dfrac1{10}$
  \item[\textbf{15.}] $35$
  \item[\textbf{16.}] ①③
\end{itemize}

\gkpart{解答题}
\begin{itemize}[leftmargin=2.6em]
  \item[\textbf{17.}] (1) $A=\dfrac\pi3$；(2) $\tan C=\dfrac{8+5\sqrt3}{11}$
  \item[\textbf{18.}] (1) $0.902$；(2) $0.254$
  \item[\textbf{19.}] (1) $MN\parallel$ 平面 $ADD_1A_1$；(2) $\arccos\dfrac2{\sqrt{21}}$；(3) $\dfrac{a^3}{6}$
  \item[\textbf{20.}] (1) $U_n=2\ln(n!)$；(2) 当 $0<x\leqslant1$ 时为 $1$；当 $x>1$ 时为 $\dfrac1{x^2}$
  \item[\textbf{21.}] $m=4$，$S=\sqrt3$
  \item[\textbf{22.}] (1) 见解析版；(2) 见解析版
\end{itemize}
